\documentclass[10pt]{article}

\usepackage[margin=0.55in]{geometry}
\usepackage[T1]{fontenc}
\usepackage[utf8]{inputenc}
\usepackage{lmodern}
\usepackage{booktabs}
\usepackage{array}
\usepackage{enumitem}
\usepackage[hidelinks]{hyperref}

\setlength{\parindent}{0pt}
\setlength{\parskip}{0.28em}
\setlist[itemize]{leftmargin=1.2em, itemsep=0.15em, topsep=0.15em}
\setlength{\tabcolsep}{4pt}
\renewcommand{\arraystretch}{1.08}

\begin{document}

\small

\begin{center}
{\Large \textbf{Project Publication Plan}}\\
\textbf{Optimal Classification of Low-volume Medical Data with Adaptive Kernel Design}\\
20 April 2026
\end{center}

\textbf{Goal.} The goal of this project is to turn the current thesis work on adaptive kernel design into a publishable medical-imaging paper for low-volume mammographic classification.

\textbf{Current position.} The current work includes kernel-bank selection, kernel composition, rotation-aware kernels, and a final signed radial kernel with learnable radius for low-volume mammographic patch classification. The main missing parts for publication are unified evaluation, stronger baselines, better statistics, and more validation.

\textbf{Target journals and what is still missing}

\begin{tabular}{>{\raggedright\arraybackslash}p{0.22\linewidth}>{\raggedright\arraybackslash}p{0.16\linewidth}>{\raggedright\arraybackslash}p{0.56\linewidth}}
\toprule
\textbf{Journal} & \textbf{Deadline} & \textbf{Missing parts before submission} \\
\midrule
Medical Image Analysis (MedIA) & Rolling & External or patient-level validation, serious deep-learning and radiomics baselines, confidence intervals, stronger clinical framing. \\
IEEE Transactions on Medical Imaging (TMI) & Rolling & Clearer methodological novelty, stricter ablations, stronger statistical comparisons, and a fully unified leakage-safe protocol. \\
Artificial Intelligence in Medicine & Rolling & Strong baselines, repeated evaluation, better discussion of clinical relevance and interpretability. \\
Computerized Medical Imaging and Graphics (CMIG) & Rolling & Unified final experiments, cleaner heatmap/localization evidence, stronger manuscript organization. \\
Computer Methods and Programs in Biomedicine (CMPB) & Rolling & Complete benchmark table, reproducible scripts, and polished practical presentation of the method. \\
\bottomrule
\end{tabular}

\textbf{Plan and goal.}
\begin{itemize}
    \item Publish the project as a medical-imaging journal paper in 2026.
    \item Use the signed learnable-radius radial kernel as the main method.
    \item Finish the missing evaluation, baselines, statistics, and final manuscript packaging before submission.
\end{itemize}

\textbf{Main tasks that must be completed}

\begin{itemize}
    \item Freeze one final train/validation/test protocol with no patient or image leakage.
    \item Re-run the strongest bank-based, composition, and radial models on exactly the same split.
    \item Add at least two strong baselines: one classical radiomics baseline and one compact deep-learning baseline.
    \item Add confidence intervals and statistical comparison for AUC and accuracy.
    \item Keep the signed learnable-radius radial model as the main proposed method and treat unconstrained radial results only as ablations.
    \item Strengthen localization and interpretability evidence with cleaner heatmaps and explanation of why the radial prior matches lesion morphology.
    \item Convert the final experiments into clean reproducible scripts and final publication-quality tables and figures.
\end{itemize}

\textbf{Timeline}

\begin{tabular}{>{\raggedright\arraybackslash}p{0.22\linewidth}>{\raggedright\arraybackslash}p{0.72\linewidth}}
\toprule
\textbf{Date} & \textbf{Task} \\
\midrule
1 May 2026 & Finish the deep-learning baseline. \\
8 May 2026 & Finalize the evaluation protocol and the exact main claim of the paper. \\
15 May 2026 & Align the baseline and all metrics with the final benchmark setup. \\
22 May 2026 & Re-run the strongest bank-based, composition, and radial models on one common split. \\
29 May 2026 & Freeze the final benchmark setup and main comparison table. \\
5 June 2026 & Compute confidence intervals for the main metrics. \\
12 June 2026 & Finish the key ablation experiments. \\
19 June 2026 & Complete seed-stability checks and error analysis. \\
26 June 2026 & Finalize localization figures and interpretability outputs. \\
3 July 2026 & Finalize reproducible scripts and publication-ready tables. \\
10 July 2026 & Write the methods and results sections in near-final form. \\
17 July 2026 & Complete the full manuscript draft. \\
24 July 2026 & Finish internal revision and figure polishing. \\
31 July 2026 & Final journal formatting, references, and submission package. \\
7 August 2026 & Submit to MedIA/TMI if the evidence is strong enough; otherwise submit to AI in Medicine or CMIG. \\
\bottomrule
\end{tabular}

\end{document}
